% Generated draft: requires mathematical and source review.
\documentclass[11pt]{article}
\usepackage[margin=25mm]{geometry}
\usepackage{fontspec}
\setmainfont{Latin Modern Roman}
\usepackage{amsmath,amssymb,mathtools,unicode-math}
\setmathfont{Latin Modern Math}
\usepackage{enumitem,needspace,xurl,hyperref}
\setlength{\parindent}{0pt}
\setlength{\parskip}{5pt}
\setlength{\emergencystretch}{3em}
\newcommand{\N}{\mathbb{N}}
\newcommand{\Z}{\mathbb{Z}}
\newcommand{\Q}{\mathbb{Q}}
\newcommand{\R}{\mathbb{R}}
\newcommand{\C}{\mathbb{C}}
\newcommand{\F}{\mathbb{F}}
\newcommand{\A}{\mathbb{A}}
\newcommand{\PP}{\mathbb{P}}
\newcommand{\E}{\mathbb{E}}
\newcommand{\eps}{\varepsilon}
\newcommand{\dd}{\,\mathrm{d}}
\begin{document}
\section*{UnsolvedMath ID: 801021}
\textbf{Source problem code:} catalogue entry
\subsection{Definitions and mathematical statement}
Write
\[
T_R=\text{2-RAN},\qquad
T_C=\operatorname{CARD}(\Sigma_2),\qquad
T_W=\operatorname{WPHP}(\Sigma_2),\qquad
T_G=\operatorname{GPHP}(\Sigma_2).
\]
These are labels for the theories or principles named in the supplied record; their axioms are not supplied. In particular, the record does not define the relevant class \(\Sigma_2\), the underlying first-order language, or the background theory.

The source describes \(T_R\) as a first-order theory and states that it is ``between'' \(T_C\) and \(T_W\). It then asks for the relationship between \(T_G\) and \(T_R\). The supplied assertion about being between two principles is retained as source context; the record does not specify its implication directions or the comparison relation.

\textbf{Problem, subject to unresolved definitions.} Determine the relative strength of \(\operatorname{GPHP}(\Sigma_2)\) and \(\text{2-RAN}\) in the setting intended by the source.

If the intended comparison is deductive strength over a common background theory \(B\), this would ask which of
\[
B+T_R\vdash T_G,
\qquad
B+T_G\vdash T_R
\]
hold, where \(B+T\) means adjoining the axioms of \(T\) to \(B\), and derivability of a theory or scheme means derivability of every axiom or instance. This is a conditional interpretation only: neither \(B\), the axiom sets, nor deductive strength as the intended comparison is specified in the supplied record. A fully determined formal statement therefore requires additional source definitions.
\subsection{Short English statement}
The source asks how \(\operatorname{GPHP}(\Sigma_2)\) compares in strength with the first-order theory \(\text{2-RAN}\), against the supplied background assertion that \(\text{2-RAN}\) lies between \(\operatorname{CARD}(\Sigma_2)\) and \(\operatorname{WPHP}(\Sigma_2)\). The definitions, background theory, and exact comparison relation are missing, so the displayed implication questions are explicitly conditional.
\subsection{Sources}
\begin{itemize}
\item Local UnsolvedMath registry: permanent ID \(801021\), source problem code \texttt{catalogue entry}, title \emph{Open Problems in Mathematical Logic --- Problem 33}; collection: AMR Open Problem Lists. The registry attributes the record to Barmpalias, Gao et al., with proposed year 2026, and labels it open. These are supplied registry attributions and status labels, not independently verified mathematical conclusions.
\item Supplied source reference: Barmpalias, Gao et al., \emph{Open Problems in Mathematical Logic} (2026), Problem 33, \url{https://arxiv.org/abs/2608.26628}. The supplied provenance records access on 2026-10-16 and extraction by \texttt{arxiv\_latex}. This access date and provenance are supplied metadata; the reference has not been independently visited or verified for this draft.
\end{itemize}
\end{document}
